\section{Signals and Predictive Models}

PML evaluates a common path-prediction and decision pipeline over several
direction families, separating the quality of a direction from the evidence
used to forecast its later behavior.

\subsection{Direction Families}

\paragraph{Random control.}
A seeded unit-norm Gaussian direction defines response thresholds and the
paired null baseline. The logged \texttt{matched\_norm\_random} entry is
duplicates this direction and is retained in the audit record.

\paragraph{Mean difference.}
For positive and contrast activation sets $\mathcal{H}^+$ and
$\mathcal{H}^-$, we use
\begin{equation}
    \mathbf{v}_{\mathrm{mean}}=
    \frac{\boldsymbol{\mu}^+-\boldsymbol{\mu}^-}
    {\|\boldsymbol{\mu}^+-\boldsymbol{\mu}^-\|_2},
\end{equation}
the multi-example analogue of activation addition
\citep{turner2023steering,rimsky-etal-2024-steering}.

\paragraph{Linear and logistic probes.}
Normalized classifier weights provide supervised discriminative directions
without nonlinear feature learning.

\paragraph{RFM/AGOP direction.}
A recursive feature machine estimates task-adapted geometry through the average
gradient outer product \citep{radhakrishnan2022mechanism}. Its leading
eigenvector provides a low-rank direction, while spectrum, concentration, and
alignment statistics become geometry features. A top-$k$ boundary study tests
whether broader supervised subspaces trade selectivity for leverage.

All directions are fitted independently per record and layer and injected at
the selected residual block during evaluation. Detailed activation extraction,
fitting hyperparameters, inference settings, and continuation scoring are
included in the released protocol.

\subsection{Predictive Evidence and Grouped Evaluation}

Static localization $L$ summarizes class separation, projection, saliency, and
direction agreement; RFM geometry $G$ summarizes spectral concentration and
alignment. These feature groups are nested with baseline belief $B$, metadata
$M$, and the low-dose response $R$, allowing the contribution of static
localization to be measured separately from the observed response.

We fit class-balanced logistic regression and random forests for target,
damage, and clean outcomes. Five-fold evaluation groups all paths from the same
record, dataset, or domain, preventing related trajectories from crossing a
train--test boundary. We report prevalence, AUROC, and average precision in the
main paper;
additional classification and calibration metrics are in the released
artifacts.

\subsection{Risk-Aware Strength Selection}

For each nonzero candidate coefficient, the selector predicts clean-effect
probability and continuous utility. With $q(\alpha)\in\{-1,+1\}$ denoting the
requested suppression or enhancement sign, measured utility is
\begin{equation}
\begin{aligned}
    u_i(\alpha)={}&q(\alpha)\Delta^T_i(\alpha)
    -\max\{0,-\Delta^N_i(\alpha)\}\\
    &-\max\{0,-\Delta^C_i(\alpha)\}.
\end{aligned}
\end{equation}
Target movement is rewarded and collateral margin decreases receive unit
penalties. The decision score is
\begin{equation}
    \widehat{u}_i(\alpha)+0.1\widehat{p}_i(\mathrm{clean}\mid\alpha)
    -0.01|\alpha|.
\end{equation}
The policy selects the highest-scoring coefficient and abstains when its score
is nonpositive. This utility is defined on teacher-forced answer-margin changes.
We evaluate it against both no intervention,
whose utility is zero by construction, and a train-tuned fixed coefficient.
